% =============================================================================
% Appendices
% =============================================================================

\appendix

% -----------------------------------------------------------------------------
\section{Permutation-Sensitivity Metric Validation}
\label{app:metric_validation}
% -----------------------------------------------------------------------------

The standard positional-dependence score ($\text{pos\_dep}$) measures the mean Pearson correlation
of a head's attention pattern across different inputs at fixed positions. A score near 1 indicates
the pattern is nearly identical across inputs---classifying the head as \emph{positional}. We
discovered this metric produces systematic false positives in deep layers of models trained on
content-rich tasks: a head can attend based on token content while \emph{appearing} to have a
stable, input-independent pattern, because by layer 2 or deeper the residual stream has already
encoded stable, content-conditioned representations.

\subsection{Validation Setup}

We evaluated both metrics on three models spanning our task range:
(1) the entropy$=$0 Markov model (4 layers, 4 heads), where attention is architecturally
bypassed and all heads should genuinely be positional;
(2) the content-matching induction model (2 layers, 4 heads), where layer-1 heads are the
known induction circuit;
(3) the Shakespeare model (6 layers, 8 heads), where layers 1--5 form the suppression cascade.

For each head, we compute both $\text{pos\_dep}$ (cross-input attention pattern correlation)
and $\text{sens}_h$ (Equation~\ref{eq:perm_sens}, 50 random vocabulary permutations).
Classification thresholds: $\text{pos\_dep} > 0.6$ $\Rightarrow$ positional;
$\text{sens}_h > 0.3$ $\Rightarrow$ content-based.

\subsection{Results by Model}

\paragraph{Entropy\,$=$\,0 Markov model.} Both metrics agree on all 16 heads
(agreement rate: 16/16 = 100\%). All heads correctly classified as positional
($\text{pos\_dep}$ range: 0.982--0.998; $\text{sens}_h$ range: 0.051--0.120).
This is the expected outcome: at zero entropy, the model computes a deterministic
lookup table via MLPs; attention heads are structurally unused and develop
input-independent (genuinely positional) patterns.

\paragraph{Content-matching induction model.} Agreement on 5/8 heads (62.5\%).
The three disagreements are all layer-1 heads (L1H0, L1H2, L1H3) classified as
\emph{positional} by $\text{pos\_dep}$ but as \emph{content-based} by permutation
sensitivity. These are precisely the heads comprising the induction circuit (mean
$\text{pos\_dep} = 0.676$, mean $\text{sens}_h = 0.404$). The misclassification
occurs because L1's residual stream already carries stable token-identity information
written by L0, causing the content-conditional L1 patterns to appear repeatable across
inputs.

\paragraph{Shakespeare model.} Agreement on only 9/48 heads (18.75\%).
The 39 disagreements are exclusively layer-1 through layer-5 heads, all classified as
positional by $\text{pos\_dep}$ but as content-based by $\text{sens}_h$. No
false-content-to-positional disagreements occur in any model. The complete per-head
breakdown for the Shakespeare model is given in Table~\ref{tab:metric_shakespeare}.

\subsection{Interpretation}

The systematic direction of disagreement (always false-positional, never false-content)
confirms the mechanism: $\text{pos\_dep}$ cannot detect content-based heads in layers
where the preceding residual stream has stabilized. The permutation test forces the issue
by actively shuffling token identities, exposing content-dependence that cross-input
correlation misses.

We adopt $\text{sens}_h$ as the primary classification metric for all deep-layer heads
going forward. Layer-0 heads are reliably classified by both metrics and are unaffected
by this correction.

\begin{table}[p]
\centering
\caption{Per-head metric comparison for the Shakespeare model (6 layers, 8 heads = 48 heads
total). \texttt{pos\_dep} threshold 0.6; \texttt{perm\_sens} threshold 0.3.
Disagree$^*$ indicates a false-positive ``positional'' classification by the old metric.}
\label{tab:metric_shakespeare}
\small
\begin{tabular}{llrrll}
\toprule
Head & Layer & pos\_dep & perm\_sens & Old class. & New class. \\
\midrule
L0H0 & 0 & 0.996 & 0.091 & positional & positional \\
L0H1 & 0 & 0.997 & 0.066 & positional & positional \\
L0H2 & 0 & 1.000 & 0.021 & positional & positional \\
L0H3 & 0 & 1.000 & 0.028 & positional & positional \\
L0H4 & 0 & 0.999 & 0.033 & positional & positional \\
L0H5 & 0 & 0.999 & 0.038 & positional & positional \\
L0H6 & 0 & 0.994 & 0.099 & positional & positional \\
L0H7 & 0 & 0.995 & 0.189 & positional & positional \\
\midrule
L1H0 & 1 & 0.759 & 0.610 & positional & \textbf{content}$^*$ \\
L1H1 & 1 & 0.744 & 0.644 & positional & \textbf{content}$^*$ \\
L1H2 & 1 & 0.695 & 0.656 & positional & \textbf{content}$^*$ \\
L1H3 & 1 & 0.789 & 0.568 & positional & \textbf{content}$^*$ \\
L1H4 & 1 & 0.716 & 0.658 & positional & \textbf{content}$^*$ \\
L1H5 & 1 & 0.771 & 0.579 & positional & \textbf{content}$^*$ \\
L1H6 & 1 & 0.566 & 0.784 & content    & content \\
L1H7 & 1 & 0.851 & 0.503 & positional & \textbf{content}$^*$ \\
\midrule
L2H0 & 2 & 0.688 & 0.679 & positional & \textbf{content}$^*$ \\
L2H1 & 2 & 0.744 & 0.591 & positional & \textbf{content}$^*$ \\
L2H2 & 2 & 0.719 & 0.613 & positional & \textbf{content}$^*$ \\
L2H3 & 2 & 0.698 & 0.615 & positional & \textbf{content}$^*$ \\
L2H4 & 2 & 0.657 & 0.680 & positional & \textbf{content}$^*$ \\
L2H5 & 2 & 0.718 & 0.616 & positional & \textbf{content}$^*$ \\
L2H6 & 2 & 0.781 & 0.529 & positional & \textbf{content}$^*$ \\
L2H7 & 2 & 0.735 & 0.614 & positional & \textbf{content}$^*$ \\
\midrule
L3H0 & 3 & 0.734 & 0.586 & positional & \textbf{content}$^*$ \\
L3H1 & 3 & 0.744 & 0.571 & positional & \textbf{content}$^*$ \\
L3H2 & 3 & 0.732 & 0.569 & positional & \textbf{content}$^*$ \\
L3H3 & 3 & 0.699 & 0.620 & positional & \textbf{content}$^*$ \\
L3H4 & 3 & 0.751 & 0.556 & positional & \textbf{content}$^*$ \\
L3H5 & 3 & 0.742 & 0.587 & positional & \textbf{content}$^*$ \\
L3H6 & 3 & 0.772 & 0.537 & positional & \textbf{content}$^*$ \\
L3H7 & 3 & 0.737 & 0.583 & positional & \textbf{content}$^*$ \\
\midrule
L4H0 & 4 & 0.820 & 0.474 & positional & \textbf{content}$^*$ \\
L4H1 & 4 & 0.811 & 0.473 & positional & \textbf{content}$^*$ \\
L4H2 & 4 & 0.795 & 0.519 & positional & \textbf{content}$^*$ \\
L4H3 & 4 & 0.806 & 0.482 & positional & \textbf{content}$^*$ \\
L4H4 & 4 & 0.804 & 0.481 & positional & \textbf{content}$^*$ \\
L4H5 & 4 & 0.840 & 0.447 & positional & \textbf{content}$^*$ \\
L4H6 & 4 & 0.812 & 0.502 & positional & \textbf{content}$^*$ \\
L4H7 & 4 & 0.811 & 0.484 & positional & \textbf{content}$^*$ \\
\midrule
L5H0 & 5 & 0.855 & 0.421 & positional & \textbf{content}$^*$ \\
L5H1 & 5 & 0.890 & 0.357 & positional & \textbf{content}$^*$ \\
L5H2 & 5 & 0.836 & 0.466 & positional & \textbf{content}$^*$ \\
L5H3 & 5 & 0.850 & 0.418 & positional & \textbf{content}$^*$ \\
L5H4 & 5 & 0.867 & 0.393 & positional & \textbf{content}$^*$ \\
L5H5 & 5 & 0.875 & 0.389 & positional & \textbf{content}$^*$ \\
L5H6 & 5 & 0.846 & 0.429 & positional & \textbf{content}$^*$ \\
L5H7 & 5 & 0.868 & 0.391 & positional & \textbf{content}$^*$ \\
\midrule
\multicolumn{4}{l}{\textit{Agreement: 9/48 (18.75\%). Disagreements: 39/48 (81.25\%)}} \\
\bottomrule
\end{tabular}
\end{table}

\clearpage

% -----------------------------------------------------------------------------
\section{Ablation Protocols}
\label{app:ablation_protocols}
% -----------------------------------------------------------------------------

This appendix documents the surgical ablation implementation in detail.
The central challenge is that MLX uses lazy, functional evaluation: tensors are immutable
and do not support in-place assignment. All ablation logic must therefore be implemented
via masked multiplication rather than direct index writes.

\subsection{Core Abstraction}

The ablation entry point is \texttt{forward\_with\_ablations()}, which accepts three optional
lists specifying which components to zero:

\begin{verbatim}
def forward_with_ablations(
    model: GPT,
    idx: mx.array,
    ablate_heads: list[tuple[int, int]] | None = None,
    ablate_attn_layers: list[int] | None = None,
    ablate_mlp_layers: list[int] | None = None,
) -> mx.array:
    ...
\end{verbatim}

Components are specified as: \texttt{ablate\_heads} --- list of (layer, head) pairs to zero
individually; \texttt{ablate\_attn\_layers} --- list of layer indices for which the entire
attention output is zeroed; \texttt{ablate\_mlp\_layers} --- list of layer indices for which
the entire MLP output is zeroed. Multiple simultaneous ablations are supported by populating
more than one list.

\subsection{Single-Head Ablation}

When a specific head $(l, h)$ is ablated, the forward pass deviates from the standard path
only at layer $l$. The core logic in \texttt{\_attn\_with\_head\_ablation()} is:

\begin{enumerate}
    \item Compute QKV projections and attention weights normally (Equation~\ref{eq:attn}).
    \item Compute value-weighted outputs \texttt{out} of shape \texttt{(B, n\_heads, T, d\_head)}.
    \item Construct a head mask: a float array of shape \texttt{(1, n\_heads, 1, 1)} with 0.0
          at position $h$ and 1.0 elsewhere.
    \item Multiply: \texttt{out = out * head\_mask}.
    \item Reshape and apply the output projection as normal.
\end{enumerate}

The multiplication preserves the contributions of all non-ablated heads to the output projection
while zeroing the ablated head's value-weighted output before it reaches $W_O$.

\begin{verbatim}
# Pseudocode: head mask construction (from analysis/ablation.py)
head_mask_list = [0.0 if h == head_idx else 1.0
                  for h in range(n_heads)]
head_mask = mx.array(head_mask_list).reshape(1, n_heads, 1, 1)
out = out * head_mask  # (B, n_heads, T, d_head)
\end{verbatim}

\subsection{Full-Layer Attention Ablation}

For a layer-level attention ablation (zeroing all heads in layer $l$):

\begin{verbatim}
if layer_idx in ablate_attn_set:
    attn_out = mx.zeros_like(x_norm)
\end{verbatim}

The attention sub-block's contribution to the residual stream is replaced by a zero tensor.
The MLP sub-block and layer norm still execute normally.

\subsection{MLP Ablation}

\begin{verbatim}
if layer_idx in ablate_mlp_set:
    mlp_out = mx.zeros_like(x_norm2)
else:
    mlp_out = block.mlp(x_norm2)
x = x + mlp_out
\end{verbatim}

\subsection{Multiple Simultaneous Ablations}

All three lists are processed in a single forward pass. The priority order is:
full-layer attention ablation takes precedence over per-head ablation for the same layer
(i.e., if layer $l$ is in both \texttt{ablate\_attn\_layers} and has entries in
\texttt{ablate\_heads}, the full-layer ablation applies). MLP ablations are processed
independently of attention ablations. This allows, for example, ablating all attention
in layers 1--5 while simultaneously ablating the layer-0 MLP in a single forward pass.

\subsection{Function Signatures from Source Files}

The primary ablation function signatures are:

From \texttt{analysis/ablation.py}:
\begin{verbatim}
def forward_with_ablations(
    model: GPT,
    idx: mx.array,
    ablate_heads: Optional[list[tuple[int, int]]] = None,
    ablate_attn_layers: Optional[list[int]] = None,
    ablate_mlp_layers: Optional[list[int]] = None,
) -> mx.array

def compute_induction_accuracy(
    model: GPT,
    dataset: InductionDataset,
    n_batches: int = 10,
    batch_size: int = 32,
    ablate_heads: Optional[list] = None,
    ablate_attn_layers: Optional[list] = None,
    ablate_mlp_layers: Optional[list] = None,
) -> float

def compute_shakespeare_loss(
    model: GPT,
    dataset: TextDataset,
    n_batches: int = 10,
    batch_size: int = 16,
    ablate_heads: Optional[list] = None,
    ablate_attn_layers: Optional[list] = None,
    ablate_mlp_layers: Optional[list] = None,
) -> float
\end{verbatim}

From \texttt{analysis/ablation\_robustness.py} (simplified variant for cross-seed study):
\begin{verbatim}
def forward_with_ablations(
    model: GPT,
    idx: mx.array,
    ablate_heads: Optional[list] = None,
) -> mx.array

def compute_induction_accuracy(
    model: GPT,
    dataset: InductionDataset,
    n_batches: int = 20,
    batch_size: int = 64,
    ablate_heads: Optional[list] = None,
) -> float

def run_ablation_for_seed(
    seed: int,
    dataset: InductionDataset,
) -> dict
\end{verbatim}

\clearpage

% -----------------------------------------------------------------------------
\section{Full Results Tables}
\label{app:full_results}
% -----------------------------------------------------------------------------

\subsection{Per-Seed Ablation Results: Content-Matching Induction}
\label{app:per_seed_ablation}

Table~\ref{tab:per_seed_full} gives complete per-head ablation results for all four seeds
of the content-matching induction model. Each cell reports accuracy after ablating the
corresponding head; the critical head for each seed (largest drop) is highlighted.

\begin{table}[h]
\centering
\caption{Per-head ablation accuracy for content-matching induction, all 4 seeds.
Baseline accuracy shown in the header row. Critical head per seed in \textbf{bold}.
Drop = baseline $-$ accuracy after ablation.}
\label{tab:per_seed_full}
\small
\begin{tabular}{lrrrrrrrr}
\toprule
 & \multicolumn{2}{c}{Seed 0 (base=0.854)} & \multicolumn{2}{c}{Seed 1 (base=0.964)} &
   \multicolumn{2}{c}{Seed 2 (base=0.965)} & \multicolumn{2}{c}{Seed 3 (base=0.968)} \\
\cmidrule(lr){2-3}\cmidrule(lr){4-5}\cmidrule(lr){6-7}\cmidrule(lr){8-9}
Head & Acc & Drop & Acc & Drop & Acc & Drop & Acc & Drop \\
\midrule
L0H0 & 0.813 & 0.041 & 0.964 & $-$0.001 & 0.957 & 0.008 & 0.932 & 0.036 \\
L0H1 & 0.799 & 0.055 & \textbf{0.118} & \textbf{0.846} & \textbf{0.028} & \textbf{0.937} & 0.958 & 0.010 \\
L0H2 & 0.789 & 0.065 & 0.204 & 0.760 & 0.956 & 0.009 & 0.945 & 0.023 \\
L0H3 & \textbf{0.041} & \textbf{0.813} & 0.965 & $-$0.002 & 0.918 & 0.047 & \textbf{0.463} & \textbf{0.505} \\
L1H0 & 0.858 & $-$0.005 & 0.964 & $<$0.001 & 0.958 & 0.007 & 0.959 & 0.010 \\
L1H1 & 0.854 & $-$0.001 & 0.964 & $-$0.001 & 0.957 & 0.008 & 0.952 & 0.016 \\
L1H2 & 0.860 & $-$0.006 & 0.961 & 0.002 & 0.958 & 0.007 & 0.955 & 0.013 \\
L1H3 & 0.855 & $-$0.001 & 0.961 & 0.003 & 0.962 & 0.003 & 0.961 & 0.008 \\
\midrule
Critical & \multicolumn{2}{c}{L0H3, drop=0.813} & \multicolumn{2}{c}{L0H1, drop=0.846} &
           \multicolumn{2}{c}{L0H1, drop=0.937} & \multicolumn{2}{c}{L0H3, drop=0.505} \\
\bottomrule
\end{tabular}
\end{table}

\noindent Key observations: (1) The critical head is always in Layer 0 (4/4 seeds). (2) Layer-1
heads are individually expendable: ablating any single L1 head causes a drop of at most 0.016,
often near zero. (3) Seed 3's critical head (L0H3) shows a partial distribution: drop of 0.505
versus 0.813--0.937 in seeds 0--2, suggesting L0H3 in seed 3 shares induction responsibility
with a secondary head. (4) Mean critical-head drop: $0.775 \pm 0.162$ (bootstrap 95\% CI
[0.590, 0.906]; $p = 0.004$).

\subsection{Entropy Continuum: Full Results}
\label{app:entropy_full}

Table~\ref{tab:entropy_full} reports the complete circuit profiling results across all 11
entropy levels. Models: 4 layers, 4 heads, 128 $d_{\text{model}}$, vocab=64, 5000 training steps.

\begin{table}[h]
\centering
\caption{Full entropy continuum results. H = measured entropy in nats; Copy\%, Supp\%,
Pos\% = fraction of 16 heads classified in each role; MLP logit = total MLP logit attribution;
Attn logit = total attention logit attribution.}
\label{tab:entropy_full}
\small
\begin{tabular}{rrrrrrrr}
\toprule
$\alpha$ & H (nats) & Copy\% & Supp\% & Pos\% & MLP logit & Attn logit & MLP/Attn \\
\midrule
0.0 & 0.000 &  0\% & 19\% & 81\% & 0.842 & 0.282 & 2.99$\times$ \\
0.1 & 0.528 & 25\% & 75\% &  0\% & 1.738 & 0.169 & 10.3$\times$ \\
0.2 & 0.597 & 25\% & 75\% &  0\% & 1.738 & 0.178 &  9.8$\times$ \\
0.3 & 0.985 & 25\% & 75\% &  0\% & 2.373 & 0.194 & 12.2$\times$ \\
0.4 & 1.509 & 25\% & 75\% &  0\% & 3.399 & 0.353 &  9.6$\times$ \\
0.5 & 1.985 & 25\% & 75\% &  0\% & 4.669 & 0.462 & 10.1$\times$ \\
0.6 & 2.406 & 25\% & 75\% &  0\% & 6.098 & 0.451 & 13.5$\times$ \\
0.7 & 2.824 & 25\% & 75\% &  0\% & 7.470 & 0.512 & 14.6$\times$ \\
0.8 & 3.274 & 25\% & 69\% &  6\% & 8.653 & 0.664 & 13.0$\times$ \\
0.9 & 3.684 & 25\% & 69\% &  6\% & 7.528 & 0.590 & 12.8$\times$ \\
1.0 & 4.110 & 25\% & 75\% &  0\% & 1.692 & 0.286 &  5.9$\times$ \\
\bottomrule
\end{tabular}
\end{table}

\noindent The binary phase transition is visible between $\alpha=0$ and $\alpha=0.1$: at zero
entropy, 81\% of heads are positional and 0\% are copy heads; at any nonzero entropy, the
architecture immediately jumps to 25\% copy + 75\% suppression. The MLP logit contribution
peaks at $\alpha=0.8$ (H$\approx$3.3 nats), then falls at $\alpha=1.0$ because the task becomes
unlearnable (uniform distribution, no predictable structure).

\subsection{Depth Scaling: Full Results}
\label{app:depth_full}

Table~\ref{tab:depth_full} gives per-depth statistics for all Shakespeare depth-scaling models.
Layer-wise suppression scores (mean suppression score per layer) are given in
Table~\ref{tab:depth_cascade}.

\begin{table}[h]
\centering
\caption{Depth scaling results. All models: 8 heads, $d_{\text{model}}=512$, $d_{\text{ff}}=2048$,
Shakespeare char-level, 3000 training steps.
N-supp = number of suppression heads; Frac-supp = suppression fraction;
MLP/Attn = ratio of total MLP to total attention logit attribution.}
\label{tab:depth_full}
\begin{tabular}{rrrrrrrr}
\toprule
$L$ & Eval loss & N-supp & Frac-supp & N-copy & N-prev & N-other & MLP/Attn \\
\midrule
2 & 1.045 &  8/16 & 50.0\% & 7 & 1 & 0 & 1.38$\times$ \\
4 & 0.824 & 23/32 & 71.9\% & 7 & 1 & 1 & 1.14$\times$ \\
6 & 0.663 & 39/48 & 81.3\% & 7 & 2 & 0 & 1.04$\times$ \\
8 & 0.572 & 55/64 & 85.9\% & 7 & 1 & 1 & 0.97$\times$ \\
\bottomrule
\end{tabular}
\end{table}

\begin{table}[h]
\centering
\caption{Layer-wise mean suppression score by model depth. Layer 0 is always the encoding
layer (score $\approx 0$). Subsequent layers form a monotonically increasing suppression
cascade in all models.}
\label{tab:depth_cascade}
\begin{tabular}{rllllllll}
\toprule
$L$ & Layer 0 & Layer 1 & Layer 2 & Layer 3 & Layer 4 & Layer 5 & Layer 6 & Layer 7 \\
\midrule
2 & 0.00 & 1.33 & — & — & — & — & — & — \\
4 & 0.00 & 0.75 & 1.54 & 2.05 & — & — & — & — \\
6 & 0.00 & 0.56 & 1.48 & 1.73 & 1.91 & 2.04 & — & — \\
8 & 0.00 & 0.60 & 1.26 & 1.65 & 1.73 & 1.75 & 1.80 & 2.03 \\
\bottomrule
\end{tabular}
\end{table}

\noindent Three invariants hold across all depths: (1) Layer 0 has zero suppression (copy+prev-token
composition invariant: always 7 copy heads + 1 previous-token head). (2) Suppression grows
monotonically from layer 1 onward. (3) The final layer converges to a suppression score near 2.0
regardless of depth. The OLS slope for suppression fraction as a function of depth is
$0.059 \pm 0.014$ per layer (95\% CI [0.023, 0.109]; $p = 0.052$; $R^2 = 0.90$).

\subsection{Per-Head Profile: Shakespeare 6-Layer Model}
\label{app:shakespeare_heads}

Table~\ref{tab:shakespeare_all_heads} provides the complete per-head profile for the primary
Shakespeare model (6 layers, 8 heads, 512 $d_{\text{model}}$) used throughout the paper.
This is the model from the depth-6 row of Table~\ref{tab:depth_full}.

\begin{table}[h]
\centering
\caption{Full per-head profile for the Shakespeare 6-layer model. Columns: role assigned
by permutation-sensitivity metric; copy score (OV diagonal dominance); suppression score;
prev-token score; attention entropy; positional score (cross-input correlation).}
\label{tab:shakespeare_all_heads}
\tiny
\begin{tabular}{llrrrrrr}
\toprule
Head & Role & Copy score & Supp.\ score & Prev-tok & Entropy & Pos-dep \\
\midrule
L0H0 & copy          &  0.93 & 0.00 & 0.135 & 2.46 & 0.976 \\
L0H1 & copy          &  3.13 & 0.00 & 0.040 & 3.53 & 0.991 \\
L0H2 & copy          &  1.09 & 0.00 & 0.049 & 3.69 & 0.994 \\
L0H3 & copy          &  1.91 & 0.00 & 0.063 & 3.39 & 0.990 \\
L0H4 & previous-token &  0.43 & 0.00 & 0.810 & 1.10 & 0.995 \\
L0H5 & copy          &  1.15 & 0.00 & 0.062 & 3.51 & 0.993 \\
L0H6 & copy          &  1.75 & 0.00 & 0.107 & 3.43 & 0.990 \\
L0H7 & copy          &  0.62 & 0.00 & 0.082 & 2.51 & 0.989 \\
\midrule
L1H0 & suppression   & $-$0.81 & 0.79 & 0.246 & 1.47 & 0.738 \\
L1H1 & suppression   & $-$0.71 & 0.76 & 0.233 & 1.27 & 0.660 \\
L1H2 & suppression   & $-$0.40 & 0.39 & 0.237 & 1.26 & 0.653 \\
L1H3 & previous-token & $-$0.65 & 0.69 & 0.302 & 1.38 & 0.705 \\
L1H4 & suppression   & $-$0.30 & 0.32 & 0.194 & 1.40 & 0.691 \\
L1H5 & positional    & $-$0.18 & 0.13 & 0.125 & 1.33 & 0.601 \\
L1H6 & suppression   & $-$0.86 & 0.88 & 0.175 & 1.77 & 0.792 \\
L1H7 & suppression   & $-$0.45 & 0.50 & 0.144 & 1.69 & 0.756 \\
\midrule
L2H0 & suppression   & $-$1.06 & 1.11 & 0.239 & 1.55 & 0.729 \\
L2H1 & suppression   & $-$1.74 & 1.78 & 0.153 & 1.69 & 0.743 \\
L2H2 & suppression   & $-$1.57 & 1.59 & 0.193 & 1.49 & 0.700 \\
L2H3 & suppression   & $-$1.40 & 1.39 & 0.198 & 1.60 & 0.725 \\
L2H4 & suppression   & $-$1.86 & 1.89 & 0.161 & 1.68 & 0.752 \\
L2H5 & suppression   & $-$0.94 & 0.94 & 0.153 & 1.65 & 0.741 \\
L2H6 & suppression   & $-$1.43 & 1.43 & 0.102 & 1.91 & 0.787 \\
L2H7 & suppression   & $-$1.70 & 1.69 & 0.178 & 1.59 & 0.717 \\
\midrule
L3H0 & suppression   & $-$1.87 & 1.89 & 0.128 & 2.08 & 0.807 \\
L3H1 & suppression   & $-$1.25 & 1.30 & 0.070 & 2.10 & 0.764 \\
L3H2 & suppression   & $-$1.61 & 1.56 & 0.168 & 1.70 & 0.740 \\
L3H3 & suppression   & $-$2.12 & 2.14 & 0.115 & 1.74 & 0.754 \\
L3H4 & suppression   & $-$1.93 & 1.88 & 0.121 & 1.98 & 0.784 \\
L3H5 & suppression   & $-$1.38 & 1.39 & 0.112 & 1.88 & 0.772 \\
L3H6 & suppression   & $-$1.75 & 1.76 & 0.097 & 1.94 & 0.778 \\
L3H7 & suppression   & $-$1.89 & 1.92 & 0.174 & 1.64 & 0.734 \\
\midrule
L4H0 & suppression   & $-$1.83 & 1.82 & 0.119 & 1.84 & 0.758 \\
L4H1 & suppression   & $-$1.67 & 1.66 & 0.104 & 2.07 & 0.803 \\
L4H2 & suppression   & $-$2.05 & 2.05 & 0.075 & 2.07 & 0.790 \\
L4H3 & suppression   & $-$2.26 & 2.25 & 0.078 & 2.13 & 0.810 \\
L4H4 & suppression   & $-$1.95 & 1.95 & 0.097 & 2.12 & 0.808 \\
L4H5 & suppression   & $-$1.93 & 1.89 & 0.093 & 2.01 & 0.780 \\
L4H6 & suppression   & $-$1.97 & 1.99 & 0.147 & 1.71 & 0.732 \\
L4H7 & suppression   & $-$1.68 & 1.66 & 0.132 & 1.89 & 0.773 \\
\midrule
L5H0 & suppression   & $-$1.97 & 1.96 & 0.079 & 2.15 & 0.796 \\
L5H1 & suppression   & $-$2.40 & 2.38 & 0.074 & 2.14 & 0.801 \\
L5H2 & suppression   & $-$1.92 & 1.86 & 0.062 & 2.41 & 0.844 \\
L5H3 & suppression   & $-$2.55 & 2.51 & 0.054 & 2.30 & 0.830 \\
L5H4 & suppression   & $-$1.70 & 1.65 & 0.086 & 1.88 & 0.751 \\
L5H5 & suppression   & $-$1.72 & 1.69 & 0.053 & 2.18 & 0.802 \\
L5H6 & suppression   & $-$1.91 & 1.87 & 0.074 & 2.06 & 0.789 \\
L5H7 & suppression   & $-$2.34 & 2.37 & 0.080 & 2.31 & 0.823 \\
\bottomrule
\end{tabular}
\end{table}

\subsection{Cross-Seed Similarity Matrix Values}
\label{app:similarity_matrix}

The cross-seed circuit similarity is computed as the Pearson correlation between the
component-level effect vectors of each seed pair. A component-level effect vector records,
for each (layer, component-type) pair, the causal effect score from activation patching
(Equation~\ref{eq:patching}). Table~\ref{tab:similarity_summary} summarizes the distribution
of the $\binom{20}{2} = 190$ pairwise similarities.

\begin{table}[h]
\centering
\caption{Cross-seed circuit similarity distribution (20 seeds, 190 pairs).}
\label{tab:similarity_summary}
\begin{tabular}{lr}
\toprule
Statistic & Value \\
\midrule
Mean & 0.985 \\
Standard deviation & 0.002 \\
Minimum & 0.977 \\
5th percentile & 0.981 \\
Median & 0.986 \\
95th percentile & 0.988 \\
Maximum & 0.990 \\
Bootstrap 95\% CI & [0.982, 0.988] \\
\bottomrule
\end{tabular}
\end{table}

\noindent The distribution is narrow: the least-similar seed pair still achieves 0.977. This
concentration confirms that the 0.985 mean is not driven by a few highly-similar pairs;
essentially all 20 seeds converge to the same circuit at the component level.

\clearpage

% -----------------------------------------------------------------------------
\section{Training Dynamics Details}
\label{app:training_details}
% -----------------------------------------------------------------------------

\subsection{Hyperparameters by Experiment}

\begin{table}[h]
\centering
\caption{Training hyperparameters by experiment.}
\label{tab:hyperparams}
\begin{tabular}{lrrrrrr}
\toprule
Experiment & Steps & Batch & LR & WD & Warmup & Ckpt freq \\
\midrule
Positional induction (20 seeds) & 2,000 & 64 & $10^{-3}$ & 0.01 & 100 & 25 \\
IOI (20 seeds) & 2,000 & 64 & $10^{-3}$ & 0.01 & 100 & 25 \\
Content-matching induction & 20,000 & 64 & $10^{-3}$ & 0.01 & 500 & 100 \\
Markov entropy continuum & 5,000 & 64 & $10^{-3}$ & 0.01 & 200 & 500 \\
Shakespeare (6L) & 10,000 & 32 & $3\times10^{-4}$ & 0.1 & 500 & 200 \\
Depth scaling (all depths) & 3,000 & 32 & $3\times10^{-4}$ & 0.1 & 200 & 500 \\
\bottomrule
\end{tabular}
\end{table}

\subsection{Learning Rate Schedule}

All experiments use the same schedule: linear warmup from 0 to the peak learning rate over
\texttt{warmup\_steps} steps, followed by cosine decay to 0 over the remaining steps:

\begin{equation}
\text{lr}(t) = \begin{cases}
\text{lr}_{\text{max}} \cdot \frac{t}{T_{\text{warm}}} & t < T_{\text{warm}} \\[4pt]
\text{lr}_{\text{max}} \cdot \frac{1}{2}\!\left(1 + \cos\!\left(\frac{\pi(t - T_{\text{warm}})}{T - T_{\text{warm}}}\right)\right) & t \geq T_{\text{warm}}
\end{cases}
\end{equation}

\subsection{Checkpoint Frequency and Retention}

Dense checkpoints are essential for studying circuit formation dynamics. For the seed-divergence
experiments (positional induction), checkpoints are saved every 25 steps (80 checkpoints over
2,000 steps per seed, $\times$\,20 seeds = 1,600 checkpoint files total). For the content-matching
grokking study, checkpoints are saved every 100 steps; 11 checkpoints spanning the transition window
(steps 10,000--15,000) are used for the formation timeline analysis. All other experiments use
sparser checkpointing (every 200--500 steps) sufficient for loss-curve monitoring but not for
detailed formation film analysis.

Checkpoint files use the \texttt{safetensors} format and are retained in full. Total checkpoint
storage across all experiments is approximately 4 GB.

\subsection{Compute Budget}

All experiments were run on a single Apple M-series system (Apple Silicon, Metal GPU backend)
using the MLX framework \cite{mlx2023}. Approximate wall-clock times:

\begin{itemize}
    \item 20 seeds $\times$ 2,000 steps (positional induction): $\approx$20 minutes
    \item 20 seeds $\times$ 2,000 steps (IOI): $\approx$25 minutes
    \item 4 seeds $\times$ 20,000 steps (content-matching induction): $\approx$90 minutes
    \item 11 models $\times$ 5,000 steps (entropy continuum): $\approx$22 minutes
    \item 1 model $\times$ 10,000 steps (Shakespeare 6L): $\approx$96 minutes
    \item 4 models $\times$ 3,000 steps (depth scaling): $\approx$32 minutes
    \item Circuit mapping, ablation, transplantation analysis: $\approx$60 minutes
\end{itemize}

Total compute: approximately 6 hours wall-clock time on a single Apple M-series device.
Training speed ranges from approximately 80 steps/second (small 800K-parameter models) to
approximately 18 steps/second (14M-parameter Shakespeare model).

\subsection{Training Loss Curves}

The most diagnostically interesting loss curve is the content-matching induction model, which
exhibits a characteristic grokking signature: loss decreases slowly (from $\approx 6.0$ to
$\approx 5.5$ nats) over steps 0--11,000 as the model learns token statistics, then drops
sharply to below 1.0 nats over steps 11,000--14,000 as the induction circuit forms. The
task-specific accuracy trace (at induction-masked positions only) is even sharper: near-chance
for 11,000 steps, then a rapid transition from $\approx$6\% to $\approx$85\% over approximately
1,000 steps.

The Shakespeare model shows conventional smooth convergence: loss decays from $\approx$4.2 nats
at initialization to $\approx$1.1 nats by step 1,000, then continues to decline smoothly to
$\approx$0.66 nats at step 3,000 (depth-6 checkpoint used for depth-scaling analysis) and
$\approx$0.11 nats at step 6,000 (checkpoint used for circuit mapping in the main paper).

All positional induction models (seeds 0--19) exhibit a sharp phase transition around steps
175--225: loss drops from near-ceiling ($\approx$3.9 nats, close to $\log(50)$) to near-zero
in fewer than 50 steps. Three early-crystallizing seeds (seeds 11, 13, 14) show transitions
at steps 25--75.
